\section{Conclusion}\label{sec:conclusion}

Fourteen weeks of work produced three kinds of result. The first is
infrastructure: a single unit and harness in which rho, BSGS, the kangaroo
and every index-calculus variant are charged the same way on the same
one-target workloads, with sealed records, exact replays and admitted
intervals; and a ledger that refuses a claim with a missing field. The
second is a set of constant-factor improvements, the largest of which---the
closed-form roots of $\Sem3$---began as one sentence in a manuscript we were
asked to check, and ended as a proved identity, an oracle family covering
prime and binary fields, two to four summands, CPU and GPU, plain, tower
and dense-index bases, and the default decomposition strategy of the
codebase. The third is a map of what does not work, which is longer than
the list of what does: SAT that does not scale, symmetry that pays at
$n=17$ and costs at $n=31$, transfers that provably cannot change
decomposition counts, fibre-aware bases whose gap to rho grows as
$\#E^{1/6}$, literature attacks that are exhaustive search.

The honest summary is the campaign of \S\ref{sec:campaign}. On the
question the survey asks---does index calculus solve one unseen
elliptic-curve target, cold, for less than rho?---the answer on every curve
we measured, at every weight of the pricing band, is no. The closest arm is
a compact Frobenius-orbit Koblitz index calculus at $4$--$18\times$
signed-Frobenius rho, prime-field index calculus falls behind as $r^{1/6}$,
and no fitted exponent of an index-calculus arm lies below rho's. The
online $13.9\times$ at $n=61$ compared a table against no table and loses
$53$--$317\times$ to preprocessing rho with equal advice; it is retired.

Measured against its own baseline, which is the criterion the programme now
works to, index calculus keeps improving: about twenty changes landed on
14--15 October, among them a genus-3 index calculus after cubic GHS
whose exponent (about $q^{1.4}$) lies below rho's $q^{3/2}$, an $18\times$ cheaper
$k=4$ ladder and a further $18\times$ cheaper $k=4$ solve, and production
free relations on Koblitz curves ($0.960\times$). The genus-3 family is
the nearest to rho: $1.15\times$ a measured negation-map rho at $q=2^{20}$
after the second round, with a further $0.62\times$ step not yet matched
against rho. The stages that now
dominate are relation yield, factor-base size and linear algebra, and the
open problems of \S\ref{sec:open} say what would move them.
